% ---- score2 standard preamble (from: python3 score2_tools.py preamble) ----
\documentclass[10pt]{extarticle}
\usepackage[a4paper,margin=1.5cm]{geometry}
\usepackage[T1]{fontenc}
\usepackage{amsmath,amssymb}
\usepackage{graphicx}
\usepackage[export]{adjustbox}
\usepackage{multicol}
\usepackage{enumitem}
\usepackage{booktabs,array,tabularx}
\setlength{\columnsep}{0.7cm}
\setlength{\parindent}{0pt}
\setlength{\parskip}{3pt}
\pagestyle{empty}
\makeatletter
% figure registry (figdefs.tex is generated by the extract step)
\newcommand{\qfigdef}[4]{\@namedef{qf@file@#1}{#2}\@namedef{qf@w@#1}{#3}\@namedef{qf@wide@#1}{#4}}
\newcommand{\qf@check}[1]{\@ifundefined{qf@file@#1}{\PackageError{score2}{Unknown figure id #1}{}}{}}
\newcommand{\qf@img}[4][]{\includegraphics[width=#2\dimexpr\@nameuse{qf@w@#4}pt\relax,max width=#3,max height=0.45\textheight,#1]{\@nameuse{qf@file@#4}}}
% Full-width questions: finalize wraps every question that contains a wide figure in \qfullbegin...\qfullend,
% so the WHOLE question (stem, figures, options) is set across the page in reading order, then the two columns
% resume. The question's head - its number, stem and first figure (\qfullsplit marks the end of the head) -
% is kept on one page, so a number is never left at a page foot with its figure overleaf; the rest may break
% across pages normally, so no big gaps are left. Without a \qfullsplit, the whole question is kept together
% if it is under 45% of a page. Never write these yourself.
\newif\ifqf@full
\newif\ifqf@boxed
\newsavebox{\qf@box}
\newcommand{\qf@open}{\global\qf@fulltrue\par\global\qf@boxedtrue
  \setbox\qf@box\vbox\bgroup\hsize\textwidth\linewidth\textwidth}
\newcommand{\qf@emit}[1]{\par\egroup\global\qf@boxedfalse
  \ifdim\dimexpr\ht\qf@box+\dp\qf@box\relax<#1\textheight\noindent\usebox{\qf@box}\par\else\unvbox\qf@box\fi}
\newcommand{\qfullsplit}{\ifqf@boxed\qf@emit{0.85}\fi}
\newcommand{\qf@close}{\ifqf@boxed\qf@emit{0.45}\fi\global\qf@fullfalse}
\newcommand{\qfullbegin}{\end{multicols}\qf@open}
\newcommand{\qfullend}{\qf@close\begin{multicols}{2}}
\newcommand{\qfullstart}{\qf@open}
\newcommand{\qfulldone}{\qf@close}
% \qfig[scale]{id}: block figure, centred, directly where it is written (size set by the tool).
% Fallback if a wide figure is not inside a full-width question: break out of the columns in place - no floats.
\newcommand{\qfig}[2][1]{\qf@check{#2}%
  \ifqf@full\qf@blk{#1}{#2}\else
  \ifnum\@nameuse{qf@wide@#2}=1 \expandafter\@firstoftwo\else\expandafter\@secondoftwo\fi
  {\end{multicols}\begin{center}\qf@img{#1}{\textwidth}{#2}\end{center}\begin{multicols}{2}}%
  {\qf@blk{#1}{#2}}\fi}
\newcommand{\qf@blk}[2]{\par\smallskip{\centering\qf@img{#1}{\linewidth}{#2}\par}\smallskip}
% \qfigi[scale]{id}: inline figure (inside option lists / tables / minipages); never breaks columns.
\newcommand{\qfigi}[2][1]{\qf@check{#2}\qf@img[valign=c]{#1}{\linewidth}{#2}}
\makeatother
% structure macros - lint reads these to check completeness and figure placement
\newcommand{\papersec}[1]{\par\bigskip{\centering\Large\bfseries PAPER-#1\par}\medskip}
\newcommand{\partsec}[2]{\par\medskip\noindent{\large\bfseries PART-#1: #2}\par\smallskip}
\newcommand{\sect}[1]{\par\smallskip\noindent\textbf{#1}\par}
\newcommand{\qn}[1]{\par\smallskip\noindent\textbf{#1.}~}
\newcommand{\paraq}[2]{\par\smallskip\noindent\textbf{Paragraph for Questions #1 and #2}\par}
\newenvironment{qoptions}{\begin{itemize}[leftmargin=1.8em,itemsep=1pt,topsep=2pt,label={}]}{\end{itemize}}
\newcommand{\opt}[2]{\item[(#1)] #2}
% figure registry (generated by score2_tools.py finalize)
\qfigdef{p1_phy_s1_q01_1}{Pre-Score (2023-24) - Test 04 (29 Oct 2023)-img/p1_phy_s1_q01_1.png}{159.2}{0}
\qfigdef{p1_phy_s1_q02_1}{Pre-Score (2023-24) - Test 04 (29 Oct 2023)-img/p1_phy_s1_q02_1.png}{71.5}{0}
\qfigdef{p1_phy_s1_q02_A}{Pre-Score (2023-24) - Test 04 (29 Oct 2023)-img/p1_phy_s1_q02_A.png}{75.4}{0}
\qfigdef{p1_phy_s1_q02_B}{Pre-Score (2023-24) - Test 04 (29 Oct 2023)-img/p1_phy_s1_q02_B.png}{75.2}{0}
\qfigdef{p1_phy_s1_q02_C}{Pre-Score (2023-24) - Test 04 (29 Oct 2023)-img/p1_phy_s1_q02_C.png}{75.6}{0}
\qfigdef{p1_phy_s1_q02_D}{Pre-Score (2023-24) - Test 04 (29 Oct 2023)-img/p1_phy_s1_q02_D.png}{75.2}{0}
\qfigdef{p1_phy_s1_q03_1}{Pre-Score (2023-24) - Test 04 (29 Oct 2023)-img/p1_phy_s1_q03_1.png}{87.2}{0}
\qfigdef{p1_phy_s1_q04_1}{Pre-Score (2023-24) - Test 04 (29 Oct 2023)-img/p1_phy_s1_q04_1.png}{136.8}{0}
\qfigdef{p1_phy_s1_q05_1}{Pre-Score (2023-24) - Test 04 (29 Oct 2023)-img/p1_phy_s1_q05_1.png}{80.8}{0}
\qfigdef{p1_phy_s1_q06_1}{Pre-Score (2023-24) - Test 04 (29 Oct 2023)-img/p1_phy_s1_q06_1.png}{96.7}{0}
\qfigdef{p1_phy_s2_q07_1}{Pre-Score (2023-24) - Test 04 (29 Oct 2023)-img/p1_phy_s2_q07_1.png}{77.4}{0}
\qfigdef{p1_phy_s2_q08_1}{Pre-Score (2023-24) - Test 04 (29 Oct 2023)-img/p1_phy_s2_q08_1.png}{224.4}{0}
\qfigdef{p1_phy_s2_q09_1}{Pre-Score (2023-24) - Test 04 (29 Oct 2023)-img/p1_phy_s2_q09_1.png}{200.6}{0}
\qfigdef{p1_phy_s2_q10_1}{Pre-Score (2023-24) - Test 04 (29 Oct 2023)-img/p1_phy_s2_q10_1.png}{170.7}{0}
\qfigdef{p1_phy_s2_q11_1}{Pre-Score (2023-24) - Test 04 (29 Oct 2023)-img/p1_phy_s2_q11_1.png}{121.7}{0}
\qfigdef{p1_phy_s2_q12_1}{Pre-Score (2023-24) - Test 04 (29 Oct 2023)-img/p1_phy_s2_q12_1.png}{324.6}{1}
\qfigdef{p1_phy_s3_q03_1}{Pre-Score (2023-24) - Test 04 (29 Oct 2023)-img/p1_phy_s3_q03_1.png}{134.7}{0}
\qfigdef{p1_phy_s3_q04_1}{Pre-Score (2023-24) - Test 04 (29 Oct 2023)-img/p1_phy_s3_q04_1.png}{142.1}{0}
\qfigdef{p1_phy_s3_q05_1}{Pre-Score (2023-24) - Test 04 (29 Oct 2023)-img/p1_phy_s3_q05_1.png}{245.2}{0}
\qfigdef{p1_phy_s3_q06_1}{Pre-Score (2023-24) - Test 04 (29 Oct 2023)-img/p1_phy_s3_q06_1.png}{143.2}{0}
\qfigdef{p1_chem_s2_q11_1}{Pre-Score (2023-24) - Test 04 (29 Oct 2023)-img/p1_chem_s2_q11_1.png}{88.2}{0}
\qfigdef{p1_chem_s2_q12_A}{Pre-Score (2023-24) - Test 04 (29 Oct 2023)-img/p1_chem_s2_q12_A.png}{65.8}{0}
\qfigdef{p1_chem_s2_q12_B}{Pre-Score (2023-24) - Test 04 (29 Oct 2023)-img/p1_chem_s2_q12_B.png}{65.6}{0}
\qfigdef{p1_chem_s2_q12_C}{Pre-Score (2023-24) - Test 04 (29 Oct 2023)-img/p1_chem_s2_q12_C.png}{67.6}{0}
\qfigdef{p1_chem_s2_q12_D}{Pre-Score (2023-24) - Test 04 (29 Oct 2023)-img/p1_chem_s2_q12_D.png}{66.8}{0}
\qfigdef{p2_phy_s1_q02_1}{Pre-Score (2023-24) - Test 04 (29 Oct 2023)-img/p2_phy_s1_q02_1.png}{149.9}{0}
\qfigdef{p2_phy_s1_q03_1}{Pre-Score (2023-24) - Test 04 (29 Oct 2023)-img/p2_phy_s1_q03_1.png}{84.9}{0}
\qfigdef{p2_phy_s1_q04_1}{Pre-Score (2023-24) - Test 04 (29 Oct 2023)-img/p2_phy_s1_q04_1.png}{78.6}{0}
\qfigdef{p2_phy_s1_q06_1}{Pre-Score (2023-24) - Test 04 (29 Oct 2023)-img/p2_phy_s1_q06_1.png}{142.2}{0}
\qfigdef{p2_phy_s2_q01_1}{Pre-Score (2023-24) - Test 04 (29 Oct 2023)-img/p2_phy_s2_q01_1.png}{166.5}{0}
\qfigdef{p2_phy_s2_q02_1}{Pre-Score (2023-24) - Test 04 (29 Oct 2023)-img/p2_phy_s2_q02_1.png}{182.2}{0}
\qfigdef{p2_phy_s2_q03_1}{Pre-Score (2023-24) - Test 04 (29 Oct 2023)-img/p2_phy_s2_q03_1.png}{90.4}{0}
\qfigdef{p2_phy_s2_q05_1}{Pre-Score (2023-24) - Test 04 (29 Oct 2023)-img/p2_phy_s2_q05_1.png}{129.3}{0}
\qfigdef{p2_phy_s2_q06_1}{Pre-Score (2023-24) - Test 04 (29 Oct 2023)-img/p2_phy_s2_q06_1.png}{118.0}{0}
\qfigdef{p2_phy_s3_q01_1}{Pre-Score (2023-24) - Test 04 (29 Oct 2023)-img/p2_phy_s3_q01_1.png}{272.0}{1}
\qfigdef{p2_phy_s3_q02_1}{Pre-Score (2023-24) - Test 04 (29 Oct 2023)-img/p2_phy_s3_q02_1.png}{115.6}{0}
\qfigdef{p2_phy_s3_q03_1}{Pre-Score (2023-24) - Test 04 (29 Oct 2023)-img/p2_phy_s3_q03_1.png}{106.0}{0}
\qfigdef{p2_phy_s3_q04_1}{Pre-Score (2023-24) - Test 04 (29 Oct 2023)-img/p2_phy_s3_q04_1.png}{234.8}{0}
% ---- body skeleton ----
% \begin{document}
% \begin{center}{\Large\bfseries JEE (Main+Advanced) Score-II --- Test NN (DD Mon YYYY)}\end{center}
% \begin{multicols}{2}
%   \papersec{1}   (only for two-paper tests)
%   \partsec{1}{PHYSICS} \sect{Section-I (i): Single Correct} \qn{1} ... \qfig{p1_phy_q01_1} ...
% \end{multicols}
% \end{document}

\begin{document}
\begin{center}{\Large\bfseries Pre-Score (2023-24) --- Test 04 (29 Oct 2023)}\end{center}
\begin{multicols}{2}
\papersec{1}
\partsec{1}{PHYSICS}
\sect{SECTION-I(i)}
\qn{1} A highly conducting ring of radius R is perpendicular to and concentric with the axis of a long solenoid as shown in fig. The ring has a narrow gap of width d in its circumference. The solenoid has cross sectional area A and a uniform internal field of magnitude B. Now beginning at $t=0$, the solenoid current is steadily increased to so that the field magnitude at any time $t$ is given by $B(t)=B+\alpha t$ where $\alpha>0$. Assuming that no charge can flow across the gap, the end of ring which has excess of positive charge and the magnitude of induced e.m.f. in the ring are respectively
\qfig{p1_phy_s1_q01_1}
\begin{qoptions}
\opt{A}{X, $A\alpha$}
\opt{B}{X, $\pi R^2\alpha$}
\opt{C}{Y, $A\alpha$}
\opt{D}{Y, $\pi R^2\alpha$}
\end{qoptions}
\qn{2} Which of the following is the correct sketch for a magnetic field line for the current configuration having an infinite wire carrying current $I_2$ placed along the axis of a coil carrying current $I_1$.
\qfig{p1_phy_s1_q02_1}
\begin{qoptions}
\opt{A}{\qfigi{p1_phy_s1_q02_A}}
\opt{B}{\qfigi{p1_phy_s1_q02_B}}
\opt{C}{\qfigi{p1_phy_s1_q02_C}}
\opt{D}{\qfigi{p1_phy_s1_q02_D}}
\end{qoptions}
\qn{3} Square non-deformable superconducting frame having a mass $m$ and self inductance $L$ is placed in a non-uniform magnetic field given by $B_x=-\alpha x$, $B_y=0$, $B_z=B_0+\alpha z$. At $t=0$, the center of the frame is at the origin and sides are parallel to $x$ axis and $y$ axis. Acceleration due to gravity is in negative $z$-direction. If the frame is released from rest, the variation of $z$ coordinate as a function of time is :-
\qfig{p1_phy_s1_q03_1}
\begin{qoptions}
\opt{A}{$z=\dfrac{mgL}{a^4\alpha^2}\sin\!\left(\dfrac{a^2\alpha}{\sqrt{mL}}t\right)$}
\opt{B}{$z=\dfrac{mgL}{a^4\alpha^2}\sin\!\left(\dfrac{a^2\alpha}{2\sqrt{mL}}t\right)$}
\opt{C}{$z=\dfrac{mgL}{a^4\alpha^2}\left(1-\cos\!\left(\dfrac{a^2\alpha}{2\sqrt{mL}}t\right)\right)$}
\opt{D}{$z=\dfrac{mgL}{a^4\alpha^2}\left(\cos\!\left(\dfrac{a^2\alpha}{\sqrt{mL}}t\right)-1\right)$}
\end{qoptions}
\qn{4} Parallel rails of length $2L$ are fixed on a horizontal plane at a distance $\ell$ from each other. Two identical ideal batteries with EMF $\varepsilon$ are connected to their ends (see figure). On the rails there is a jumper with mass $m$, which can progressively slide along them. The whole system is placed in a uniform vertical magnetic field with induction $B$. Assuming that the jumper resistance is $R$, and the resistance per unit length of each of the rails is equal to $\rho$, the period of small oscillations that occur when the jumper is displaced from the equilibrium position.
\qfig{p1_phy_s1_q04_1}
\begin{qoptions}
\opt{A}{$2\pi\sqrt{\dfrac{mL}{\varepsilon\ell B}}$}
\opt{B}{$2\pi\sqrt{\dfrac{mL(\rho L+R)}{\varepsilon\ell B}}$}
\opt{C}{$2\pi\sqrt{\dfrac{2mL(\rho L+R)}{\varepsilon\ell B}}$}
\opt{D}{None of the above}
\end{qoptions}
\qn{5} The current-time ($I$--$t$) graph for triangular wave having peak value $I_0$ is as shown in figure. The rms value of current is :
\qfig{p1_phy_s1_q05_1}
\begin{qoptions}
\opt{A}{$I_0$}
\opt{B}{$\sqrt{\dfrac23}\,I_0$}
\opt{C}{$\dfrac{I_0}{\sqrt2}$}
\opt{D}{$\dfrac{I_0}{\sqrt3}$}
\end{qoptions}
\qn{6} The uniform rod has a mass of $m$. If it is acted upon by a periodic force of $F=F_0\sin\omega t$, determine the angular amplitude of the steady-state vibration.
\qfig{p1_phy_s1_q06_1}
\begin{qoptions}
\opt{A}{$\dfrac{F_0}{mg+kL}$}
\opt{B}{$\dfrac{F_0}{2(mg+kL)-m\omega^2L}$}
\opt{C}{$\dfrac{3F_0}{\frac32(mg+kL)-m\omega^2L}$}
\opt{D}{$\dfrac{F_0}{3(mg+kL)-2m\omega^2L}$}
\end{qoptions}
\sect{SECTION-I(ii)}
\qn{7} A metallic pulley is in the shape of a disc of radius $a$. It can rotate freely about a horizontal axis passing through its centre. The moment of inertia of the pulley about this axis is $I$. A light string is tightly wrapped around the pulley with its one end connected to a block of mass $M$. The centre of the pulley and its circumference are connected to a resistance $R$ as shown. The contact of resistance at the circumference does not cause any friction. A uniform magnetic field $B$ is switched on which is parallel to the axis of rotation of the pulley (see Figure). The mass $M$ is allowed to fall. Assume that resistivity of the material of the pulley is negligible.
\qfig{p1_phy_s2_q07_1}
\begin{qoptions}
\opt{A}{The acceleration of the block of mass $M$ at the instant its velocity becomes $v_0$ is $\dfrac{Mga^2-\dfrac{B^2a^4v_0}{4R}}{I+Ma^2}$.}
\opt{B}{Assuming that the block can fall through a large distance, the terminal speed ($v_T$) that it will acquire is $\dfrac{4MgR}{B^2a^2}$.}
\opt{C}{The rate of change of kinetic energy of the pulley at the instant when the speed of the falling block is $\dfrac{v_T}{2}$ is equal to $\dfrac{IM^2g^2R}{B^2a^2(I+Ma^2)}$.}
\opt{D}{The rate of change of kinetic energy of the pulley at the instant when the speed of the falling block is $\dfrac{v_T}{2}$ is equal to $\dfrac{IM^2g^2R}{2B^2a^2(I+Ma^2)}$.}
\end{qoptions}
\qn{8} Cylindrical magnet of mass $m$ and magnetic moment $p_m$, is attached to the spring of stiffness $k$ such that it can oscillate along the horizontal axis which is directed along the magnetic moment. At some distance $z$ from the equilibrium position of the magnet a small metal disc is placed such that its axis coincides with the axis of the magnet. The disk has radius $R$ and thickness $h$ ($h\ll R\ll z$), the electrical resistivity of the disk material is $\rho$, and the magnetic permeability is equal to $\mu=1$. The magnet is moved from the equilibrium position and starts performing small oscillations. When the speed of magnet is $v$,
\qfig{p1_phy_s2_q08_1}
\begin{qoptions}
\opt{A}{The force exerted by disc on the magnet is equal to $\dfrac{9}{32}\dfrac{\mu_0^2p_m^2}{\pi\rho z^8}hR^4v$.}
\opt{B}{The angular frequency of oscillation is $\sqrt{\dfrac km}$.}
\opt{C}{The rate of power loss in the disk is $\dfrac{9h\mu_0^2p_m^2R^4}{32\pi z^8\rho}v^2$.}
\opt{D}{The rate of power loss in the disk is $\dfrac{9h\mu_0^2p_m^2R^4}{16\pi z^8\rho}v^2$.}
\end{qoptions}
\qn{9} A wire mesh having identical resistance $r$ as shown in the figure-1 is placed in a uniform magnetic field in the direction perpendicular to the plane of the mesh. The magnetic field is changing with time such that the rate of change of magnetic flux through one square is $\left|\dfrac{d\phi}{dt}\right|=4\varepsilon_0$. The resistance in wire mesh as shown in figure-1 are replaced by identical capacitance $C$ as shown in figure-2.
\qfig{p1_phy_s2_q09_1}
\begin{qoptions}
\opt{A}{The current in the resistance XY in circuit-1 is $\dfrac{\varepsilon_0}{r}$.}
\opt{B}{The current in the resistance LM in circuit-1 is $\dfrac{7\varepsilon_0}{2r}$.}
\opt{C}{The current in the resistance XY in circuit-2 is $\dfrac{\varepsilon_0}{4r}$.}
\opt{D}{The charge on the capacitance LM in circuit-2 is $5C\varepsilon_0$.}
\end{qoptions}
\qn{10} A steady current 4 A flows in an inductor coil when connected to a 12V dc source as shown in figure-1. If the same coil is connected to an ac source of 12V, 50 rad/s, a current of 2.4 A flows in the circuit as shown in figure-2. Now after these observations, a capacitor of capacitance $\dfrac1{50}$F is connected in series with the coil and with the same AC source as shown in figure-3.
\qfig{p1_phy_s2_q10_1}
\begin{qoptions}
\opt{A}{The inductance of the coil is nearly 0.02 H}
\opt{B}{The inductance of the coil is nearly 0.04 H}
\opt{C}{The average power supplied to the circuit after connecting capacitor in series is 30 W.}
\opt{D}{The average power supplied to the circuit after connecting capacitor in series is 24 W.}
\end{qoptions}
\qn{11} In the shown circuit $R_1=10\,\Omega$, $L=\dfrac{\sqrt3}{10}$H, $R_2=20\,\Omega$, $C=\dfrac{\sqrt3}{2}$mF. At the instant when current through $R_1$ is $10\sqrt2$A,
\qfig{p1_phy_s2_q11_1}
\begin{qoptions}
\opt{A}{The current through $R_2$ is $5\sqrt6$A.}
\opt{B}{The current through $R_2$ is zero.}
\opt{C}{The voltage drop across the capacitance is $200\sqrt2$V}
\opt{D}{The voltage drop across the inductance is $100\sqrt6$V.}
\end{qoptions}
\qfullbegin
\qn{12} Four long wires, with identical currents, either directly into or directly out of the page, are placed at four corners of a squares (of side $a$) in four different arrangement as shown. $B_1$, $B_2$, $B_3$ and $B_4$ are the magnitude of net magnetic field at the center of the square in respective cases.
\qfig{p1_phy_s2_q12_1}
\qfullsplit
\begin{qoptions}
\opt{A}{$B_1=B_4=0$}
\opt{B}{$B_2=\dfrac{\sqrt2\mu_0 i}{\pi a}$}
\opt{C}{$B_3=\dfrac{2\mu_0 i}{\pi a}$}
\opt{D}{$B_2=B_3$}
\end{qoptions}
\qfullend
\sect{SECTION-II}
\qn{1} A deflection of 24 divisions on a ballistic galvanometer (for ballistic galvanometer charge $\propto$ deflection) is obtained by discharging through it a capacitor of $3\,\mu$F charged to a potential difference of 2 V. When the same galvanometer is connected to a coil of 80 turns, each of diameter 1 cm, and the coil is quickly introduced into a strong magnetic field, a deflection of 16 divisions is observed. Calculate the strength of the magnetic field. The total resistance of the galvanometer circuit is 2000 $\Omega$.
\qn{2} In order to impart an angular velocity to an earth satellite, the geomagnetic field can be used. Find the maximum possible angular velocity (in $10^{-2}$ rad/s) about its own axis gained by the satellite if a storage battery with a capacity of $Q=5$ Amp--hours is discharged suddenly through a coil of $N=20$ turns wound around the satellite's surface along the circumference of the largest circle. The satellite has a mass of $m=10^3$ kg and is a thin walled uniform sphere. The geomagnetic field is parallel to the winding plane and its flux density is $B=0.5$ Gauss.
\qn{3} In the given circuit, initial charge on the capacitor is 1.5 C with polarity as shown in figure and current in the inductor is zero. Now at $t=0$ if the key k is closed, then the minimum current through the key is :-
\qfig{p1_phy_s3_q03_1}
\qn{4} A wire of mass $m$ is suspended from two points on the same horizontal level as shown in the diagram. The wire is placed in a uniform horizontal magnetic field and a current $I$ flows through it. The tension (in N) at the lowest point of the wire is ($m=1$ kg, $I=2$A, $L=2$m, $B=3$T, $\alpha=37^\circ$, $H=3$m)
\qfig{p1_phy_s3_q04_1}
\qn{5} The graph below shows the variation of power as a function of frequency of the AC source in a series RLC circuit. The quality factor for the circuit is.
\qfig{p1_phy_s3_q05_1}
\qn{6} Professor Proton has discovered a new planet on one of his planetary expeditions. He wants to measure the magnetic field of the planet he has found. Professor Proton has brought all the necessary equipment required to carry out the following experiment. A proton is launched off a large cliff at a speed $v$ and an angle $\theta=30^\circ$ with respect to the horizontal plane at the magnetic equator of the planet. The magnetic field acting perpendicularly on the particle can be assumed to be perfectly horizontal and coming out of the page, as shown in the diagram. How strong is the magnetic field at the magnetic equator of this planet if the time period of oscillation of horizontal component of velocity of the particle is $5\times10^{-5}$ s ? Write your answer in terms of $\mu$T (micro-Teslas). Take mass of proton $=1.6\times10^{-27}$ kg, charge on proton $=1.6\times10^{-19}$ C and $\pi=3.14$.
\qfig{p1_phy_s3_q06_1}
\partsec{2}{CHEMISTRY}
\sect{SECTION-I(i)}
\qn{1} At what minimum concentration of OH$^-$ ions will $10^{-3}$ mole of Zn(OH)$_2$ go into solution as Zn(OH)$_4^{2-}$ in 1 L aqueous solution.\par
Zn(OH)$_2$(s) + 2OH$^-$(aq.) $\rightleftharpoons$ Zn(OH)$_4^{2-}$(aq.) \ ($K_C=10^{-2}$)
\begin{qoptions}
\opt{A}{0.1 M}
\opt{B}{0.632 M}
\opt{C}{0.0316 M}
\opt{D}{0.316 M}
\end{qoptions}
\qn{2} A aqueous solution is a mixture of 0.06 M KCl and 0.06 M KI. Solid AgNO$_3$ is being added gradually till AgCl starts precipitating ($K_{sp}$ [AgCl] $=1\times10^{-10}$ and $K_{sp}$ [AgI] $=4\times10^{-16}$). The concentration of iodide ion at this instant will be nearly equal to :
\begin{qoptions}
\opt{A}{$4.0\times10^{-5}$ M}
\opt{B}{$4\times10^{-8}$ M}
\opt{C}{$2.4\times10^{-8}$ M}
\opt{D}{$2.4\times10^{-7}$ M}
\end{qoptions}
\qn{3} The degree of dissociation of water in a 0.1 M aqueous solution of HCl at a certain temperature $t^\circ$C is $3.6\times10^{-15}$. The temperature `$t$' must be : [Take density of water at $t^\circ$C $=1$ g/mL]
\begin{qoptions}
\opt{A}{$<25^\circ$C}
\opt{B}{$=25^\circ$C}
\opt{C}{$>25^\circ$C}
\opt{D}{Insufficient data to predict}
\end{qoptions}
\qn{4} Two moles of an ideal gas ($\gamma=3/2$) are made to expand reversibly and adiabatically to 4 times its initial volume. The change in entropy of the system during expansion is : (Given : $R=2$ cal/K/mole, $\log_{10}2=0.3$, $\log_{10}3=0.48$)
\begin{qoptions}
\opt{A}{5.6 cal/k}
\opt{B}{11.2 cal/K}
\opt{C}{2.8 cal/K}
\opt{D}{None of these}
\end{qoptions}
\qn{5} 0.5 mole each of two ideal gas A ($C_V=R$) and B ($C_V=R$) are taken in a container and expanded reversibly and adiabatically from $V=1$ litre to $V=2$ litre starting from initial temperature $T=300$ K. $\Delta H$ for the process (in cal/mol) is :
\begin{qoptions}
\opt{A}{$-500$ cal}
\opt{B}{$-900$ cal}
\opt{C}{$-450$ cal}
\opt{D}{$-600$ cal}
\end{qoptions}
\qn{6} Calculate $\Delta G$ for the given phase transition :\par
H$_2$O ($\ell$, 2 mole, 1 atm, 100$^\circ$C) $\longrightarrow$ H$_2$O (g, 2 mole, 2 atm, 100$^\circ$C)\par
Assume H$_2$O(g) to behave as an ideal gas. Give your answer in Kcal rounded off to nearest integer.
\begin{qoptions}
\opt{A}{3}
\opt{B}{2}
\opt{C}{1}
\opt{D}{4}
\end{qoptions}
\sect{SECTION-I(ii)}
\qn{7} Select correct statements :
\begin{qoptions}
\opt{A}{For an ideal gas (undergoing no association or dissociation) adiabatic free expansion \& isothermal free expansion have equal change in enthalpy for same volume change $V_1$ to $V_2$ starting from same initial state.}
\opt{B}{An ideal mono atomic gas follows the process $P^5V^7=$ constant. If heat is withdrawn from the gas during this process, then temperature of gas must increase.}
\opt{C}{For any cyclic process : $\displaystyle\oint\dfrac{dq}{T}\leqslant0$}
\opt{D}{Adiabatic process is always an isentropic process.}
\end{qoptions}
\qn{8} Which of the following is/are incorrect statement(s) about $10^{-3}$ M aqueous solution of HCl at 298 K ?
\begin{qoptions}
\opt{A}{On dilution pH will decrease}
\opt{B}{[H$_3$O$^+$] from H$_2$O is nearly $10^{-3}$ M}
\opt{C}{$\alpha$ for H$_2$O is nearly $1.8\times10^{-13}$}
\opt{D}{[OH$^-$] in solution is nearly $10^{-12}$}
\end{qoptions}
\qn{9} A 40.0 ml solution of weak base, BOH is titrated with 0.1 N HCl solution. The pH of the solution is found to be 10.0 and 9.0 after adding 5.0 ml and 20.0 ml of the acid, respectively. The correct informations is/are : ($\log 2=0.3$)
\begin{qoptions}
\opt{A}{$K_b=2\times10^{-5}$}
\opt{B}{$pK_b=4.5$}
\opt{C}{Initially : [BOH] $=\dfrac{3}{40}$M}
\opt{D}{pH at equivalence point $<7$.}
\end{qoptions}
\qn{10} Choose the incorrect statement(s) among the following options :
\begin{qoptions}
\opt{A}{Change in Gibb's function at constant temperature and Pressure in any process involving no non-PV work is negative.}
\opt{B}{For a reaction whose standard entropy change is positive, as temperature is increased, equilibrium constant for the reaction increases necessarily.}
\opt{C}{Exothermic reaction can not be spontaneous at higher temperature.}
\opt{D}{Magnitude of change in Gibb's function at constant temperature and pressure is equal to magnitude of maximum non-PV work that can be done by the system.}
\end{qoptions}
\qn{11} Consider an isothermal cylinder massless piston assembly in which ideal gas is filled. Cross sectional area of the cylinder $=1\,\text{m}^2$. The masses $m_1$, $m_2$ and $m_3$ are kept on the piston. When $m_1$ is removed, piston moves upto point A. When $m_1$ \& $m_2$ both are removed piston moves upto point B and when $m_1$, $m_2$ and $m_3$ all three are removed, piston moves upto point C.\par
[Given : $m_1=2\times10^4$ kg, $m_2=3\times10^4$ kg, $g=10$ m/s$^2$]
\qfig{p1_chem_s2_q11_1}
\begin{qoptions}
\opt{A}{`$m_3$' $=5\times10^4$ kg}
\opt{B}{Volume of gas at point B is 10 litre}
\opt{C}{Pressure of gas at point B is $5\times10^5$ Pa}
\opt{D}{Work done when piston moves from B to C is $-3200$ J}
\end{qoptions}
\qn{12} An ideal gas undergoes a reversible isothermal expansion from state I to state II followed by a reversible adiabatic expansion from state II to state III. The correct plot(s) representing the changes from state I to state III is(are)\par
(p : pressure, V : volume, T : temperature, H: enthalpy, S : entropy)
\begin{qoptions}
\opt{A}{\qfigi{p1_chem_s2_q12_A}}
\opt{B}{\qfigi{p1_chem_s2_q12_B}}
\opt{C}{\qfigi{p1_chem_s2_q12_C}}
\opt{D}{\qfigi{p1_chem_s2_q12_D}}
\end{qoptions}
\sect{SECTION-II}
\qn{1} Calculate pH at equilibrium of the solution in which following reaction taking place under basic medium at 300 K.\par
I$_2$(s) $\leftrightarrow$ I$^-$(aq) + IO$_3^-$(aq)\par
Given
\begin{center}
\footnotesize\setlength{\tabcolsep}{2pt}\renewcommand{\arraystretch}{1.4}
\begin{tabular}{|l|c|c|c|c|}\hline
Subst. & I$^-$(aq) & IO$_3^-$(aq) & H$_2$O($\ell$) & OH$^-$(aq)\\\hline
$\Delta G^\circ_f$ (kJ mol$^{-1}$) & $-50$ & $-123.5$ & $-233$ & $-150$\\\hline
\end{tabular}
\end{center}
Equilibrium concentration of I$^-$ $=0.1$ M\par
Equilibrium concentration of IO$_3^-$ $=0.1$ M\par
$R=\dfrac{25}{3}$ J mol$^{-1}$ K$^{-1}$; $\ell$n 10 $=2.3$
\qn{2} The vapour pressures of liquid water and ice at $-10^\circ$C are 0.32 and 0.16 Pa, respectively. What is the magnitude of Gibbs free energy change for the process (in Joule).\par
H$_2$O($\ell$, $-10^\circ$C, 0.32Pa, 1 mole) $\longrightarrow$ H$_2$O (s, $-10^\circ$C, 0.16 Pa, 1 mole)\par
[Given : $R=8.3$ J K$^{-1}$ mol$^{-1}$, $\ell$n 2 $=0.7$]\par
(Answer to the nearest integer)
\qn{3} 2A(s) $\rightleftharpoons$ B(g) + 2C(g) + 3D(g)\par
Total pressure developed in a closed container by decomposition of A at equilibrium is 12 bar at 727$^\circ$C. Calculate $|\Delta G^\circ|$ (in kcal) of the reaction at 727$^\circ$C ($R=2$ cal/mole-K, ln 2 = 0.7, ln 3 = 1.1)
\qn{4} When 1000 bar constant external pressure is applied on FCC crystal of metal `M' then its crystalline structure is changed from FCC to BCC. If radius of metal is $100\sqrt3$\,\AA\ and initially in the FCC lattice there were $10^{15}$ unit cells, then calculate work done by the metal crystal, in joule. [Take $\sqrt6=2.5$]
\qn{5} O$_2$ is kept in a vessel at 300 K. Following equilibrium is established $3\text{O}_2(g)\rightleftharpoons2\text{O}_3(g)$. At equilibrium partial pressure of O$_2$ $=10^5$ bar and partial pressure of O$_3$ $=10^{-10}$ bar. Calculate $\Delta_rG^\circ$ (in Kcal) of the above reaction.\par
[Given : $R=2$ cal K$^{-1}$ mol$^{-1}$; log$_e$ 10 $=2.3$]
\qn{6} 500 ml of an aqueous solution of 0.01 M AgNO$_3$ is mixed with 250 ml of aqueous solutions of each of NaBr and NaCl, each having molarity 0.02 M. The equilibrium concentration of Br$^-$ when AgCl starts precipitating is $x\times10^{-y}$ mole/L. The value of `x' is :\par
[Given $K_{sp}$ (AgBr) $=5\times10^{-13}$; $K_{sp}$ (AgCl) $=10^{-10}$].
\partsec{3}{MATHEMATICS}
\sect{SECTION-I(i)}
\qn{1} Two triangles ABC and DEF are formed by three tangents and the corresponding three normals to a parabola $y^2=4ax$ respectively. The orthocenter of $\Delta$ABC has coordinates $(-5,4)$. Then, distance of orthocenter of $\Delta$DEF from axis of parabola is
\begin{qoptions}
\opt{A}{5 units}
\opt{B}{2 units}
\opt{C}{4 units}
\opt{D}{cannot be determined}
\end{qoptions}
\qn{2} Let P be any point on a directrix of an ellipse of eccentricity $e$, S be the corresponding focus and C is the centre of the ellipse. The line PC meets the ellipse at A and the angle between PS and tangent at A is $\alpha$, then $\alpha$ is equal to -
\begin{qoptions}
\opt{A}{$\tan^{-1}e$}
\opt{B}{$\dfrac{\pi}{2}$}
\opt{C}{$\tan(1-e^2)$}
\opt{D}{none of these}
\end{qoptions}
\qn{3} For an ellipse $\dfrac{x^2}{9}+\dfrac{y^2}{4}=1$, with vertices A \& A$'$, tangent drawn at the point `P' in the first quadrant meets the $y$-axis in Q and the chord A$'$P meets the $y$-axis in M. If `O' is the centre of ellipse, then $OQ^2-MQ^2$ is equal to
\begin{qoptions}
\opt{A}{9}
\opt{B}{13}
\opt{C}{4}
\opt{D}{5}
\end{qoptions}
\qn{4} The normal at P to a hyperbola of eccentricity $e$, intersects its transverse and conjugate axes at L and M respectively. If locus of the mid-point of LM is hyperbola, then eccentricity of the hyperbola is :
\begin{qoptions}
\opt{A}{$\left(\dfrac{e+1}{e-1}\right)$}
\opt{B}{$\dfrac{e}{\sqrt{(e^2-1)}}$}
\opt{C}{$e$}
\opt{D}{None of these}
\end{qoptions}
\qn{5} Let AB and CD are two chords of hyperbola $9x^2-y^2=1$. If AB and CD are perpendicular and S(1,0) is their point of intersection, then the value of $\dfrac{1}{(AS)(BS)}+\dfrac{1}{(CS)(DS)}$ is equal to
\begin{qoptions}
\opt{A}{1}
\opt{B}{8}
\opt{C}{9}
\opt{D}{$\dfrac19$}
\end{qoptions}
\qn{6} The locus of point of intersection of tangents drawn at ends of all normal chords to parabola $y^2=8x-8$ is -
\begin{qoptions}
\opt{A}{$y^2(x-3)=-32$}
\opt{B}{$y^2(x+3)=-32$}
\opt{C}{$x^2(y+3)=-32$}
\opt{D}{$x^2(y-3)=-32$}
\end{qoptions}
\sect{SECTION-I(ii)}
\qn{7} A tangent is drawn at any point P($t$) on the parabola $y^2=8x$ and a point Q$(\alpha,\beta)$ is taken on it from which pair of tangents QA and QB are drawn to the circle $x^2+y^2=4$, where A \& B are points of contact. Now if point P($t$) varies on the parabola then the locus of all such point of concurrency of the chord of contact AB is denoted by $S=0$, then -
\begin{qoptions}
\opt{A}{Equation of $S=0$ is $2x+y^2=0$}
\opt{B}{Equation of $S=0$ is $x+y^2=0$}
\opt{C}{Area bounded by $S=0$ and the line $2x+1=0$ is $\dfrac23$ sq. units}
\opt{D}{Tangents drawn at the extremities of intersection of $S=0$ and the line $2x-y+1=0$ to the curve $S=0$ includes an angle $75^\circ$.}
\end{qoptions}
\qn{8} Let P$(x_0,y_0)$ be any point on the parabola $y^2=4x$, if normal at P intersect the parabola again at Q and tangents at P and Q intersect at point R -
\begin{qoptions}
\opt{A}{acute angle between normal at P and tangent at Q to the parabola is $\theta$, then $\tan\theta=\dfrac{|y_0|}{4}$}
\opt{B}{length of normal chord PQ is $\dfrac{4(x_0+1)^{3/2}}{x_0}$}
\opt{C}{ordinate of point R is $-\dfrac{4}{y_0}$}
\opt{D}{ordinate of circumcentre of triangle PQR is $\dfrac{-2(x_0+3)}{y_0}$}
\end{qoptions}
\qn{9} Let $S_1$ is a hyperbola and $S_2$ be an ellipse whose centers are at origin. The transverse axis of $S_1$ and major axis of $S_2$ lie along the $x$-axis and $y$-axis respectively. Let H be the hyperbola $x^2-y^2=7$. The straight line $4x-3y=7$ touches the curves H, $S_1$ and $S_2$ at P, Q and R respectively. Suppose PQ $=\dfrac52=$ QR. If $e_1$ and $e_2$ are eccentricities of $S_1$ and $S_2$ respectively. Then -
\begin{qoptions}
\opt{A}{$e_1=\dfrac{23}{15}$}
\opt{B}{$e_2=\dfrac12$}
\opt{C}{$15e_1^2+4e_2^2=24$}
\opt{D}{$e_1$ is an irrational number}
\end{qoptions}
\qn{10} Let H : $y(3y+4x)=-4$ is a hyperbola and $y=mx+c$ is its conjugate axis. Length of latus rectum of H is $L$, eccentricity $e$ and $(x_1,y_1)$ is one vertex with $y_1>0$, then -
\begin{qoptions}
\opt{A}{$m=2$}
\opt{B}{$L=1$}
\opt{C}{$4e^2=5$}
\opt{D}{$5y_1^2=4$}
\end{qoptions}
\qn{11} Let two parabolas have a common axis where focus of each being exterior to the other and let $\ell_1$, $\ell_2$ be their latus rectums then the locus of the mid points of the intercepts between the parabolas made on the lines parallel to the common axis is a
\begin{qoptions}
\opt{A}{Straight line if $\ell_1=\ell_2\ \forall\ \ell_1,\ell_2\in R^+$}
\opt{B}{Parabola if $\ell_1\neq\ell_2\ \forall\ \ell_1,\ell_2\in R^+$}
\opt{C}{Parabola $\forall\ \ell_1,\ell_2\in R^+$}
\opt{D}{Hyperbola if $\ell_1=2\ell_2\ \forall\ \ell_1,\ell_2\in R^+$}
\end{qoptions}
\qn{12} A parabola whose focus is S(6, 8) is touching the co-ordinate axes, then which of the following statement(s) is/are correct ?
\begin{qoptions}
\opt{A}{The equation of circle whose diameter is the portion of the tangent at the vertex of the parabola between the co-ordinate axis is $x^2+y^2-6x-8y=0$.}
\opt{B}{If P and Q are the ends of a focal chord of the parabola, then $\dfrac1{SP}+\dfrac1{SQ}=\dfrac5{24}$.}
\opt{C}{Equation of the latus rectum is $4x+3y=48$.}
\opt{D}{If P and Q are the ends of a focal chord of the parabola, then $\dfrac1{SP}+\dfrac1{SQ}=\dfrac15$.}
\end{qoptions}
\sect{SECTION-II}
\qn{1} Chord of the curve $3x^2-y^2-2x+4y=0$ which subtends a right angle at the origin, always passes through vertex of a concave up parabola whose axis is parallel to $y$-axis and length of latus rectum is 8. If extremity of latus rectum of parabola is $L_1(\alpha,\beta)$ and $L_2(\gamma,\delta)$, then find the value of $(\alpha+\beta+\gamma+\delta)$ ?
\qn{2} One of the focus of ellipse $\dfrac{x^2}{8}+\dfrac{y^2}{4}=1$ is at S and P is the upper end of its minor axis. A parabola is drawn whose vertex is at origin and P, S are extremities of its latus rectum, then one of its possible equations is of the form $\lambda(x-1)^2+\lambda\{y-1\}^2=(x+y+\lambda)^2$, where $\lambda$ is
\qn{3} From a point P on parabola $y^2=4x$ a normal is drawn to intersect it at Q and tangent drawn at P intersect $x$-axis at R. If the locus of centroid of $\Delta$PQR is $108xy^2=(a+by^2)^2$, then the value of $(b-a)$ is
\qn{4} Consider an ellipse E : $\dfrac{x^2}{a^2}+\dfrac{y^2}{b^2}=1$ and a parabola P : $y^2=4c(x+a)$, ($a>0$). A straight line intersects `E' at points A and B and `P' at points C and D. If the angle subtended by AB at the centre of the ellipse is $90^\circ$ then the locus of point of intersection of tangents to the parabola `P' at C and D is $a^2b^2\,(y^2+\alpha c^2)=4c^2(x+\beta a)^2\,(a^2+b^2)$. The value of $\alpha$ is
\qn{5} Let $(-1,-1)$ be the focus of a parabola and $3x-y=8$ be the tangent drawn to it at point $(7,13)$ on the parabola. If $ax+by+19=0$ is the equation of the directrix of the parabola then the value of $a+b$ is (where $a,b\in N$)
\qn{6} Consider an ellipse $4x^2=52-13y^2$ and a variable point P on the line $y+2x=12$ such that $\angle F_1PF_2$ is maximum where $F_1$ and $F_2$ are the foci of the given ellipse then $\left(\dfrac{F_2P}{F_1P}\right)^2$ is equal to (where $F_1$ lies on positive $x$-axis)
\papersec{2}
\partsec{1}{PHYSICS}
\sect{SECTION-I}
\qn{1} In an AC circuit, we have an inductance $L$ joined in series to two resistors of resistance $R$ each. In the first case, the resistors are in series with each other and in the second case, the resistors are parallel to each other. The power consumed by the circuit in the first case is 810W and in the second case, it is 506.25 W. If the source voltage is 225 V,
\begin{qoptions}
\opt{A}{The value of $R$ is 20 $\Omega$.}
\opt{B}{In the first case, the phase difference between current through the source and voltage across it is $37^\circ$.}
\opt{C}{In the second case, the phase difference between current through the source and voltage across it is $\tan^{-1}(3)$}
\opt{D}{The value of Inductive reactance is 30 $\Omega$.}
\end{qoptions}
\qn{2} A wooden cubical block ABCDEFG of mass $m$ and side $a$ is wrapped by square wire loop of perimeter $4a$, carrying current $I$. The whole system is placed at frictionless horizontal surface in a uniform magnetic field $\vec B=B_0\hat j$ as shown in figure. In this situation, normal force between horizontal surface and block passes through a point at a distance $x$ from center. Choose CORRECT statement(s) :
\qfig{p2_phy_s1_q02_1}
\begin{qoptions}
\opt{A}{The block must not topple if $I<\dfrac{mg}{aB_0}$}
\opt{B}{The block must not topple if $I<\dfrac{mg}{2aB_0}$}
\opt{C}{$x=\dfrac a4$ if $I=\dfrac{mg}{2aB_0}$}
\opt{D}{$x=\dfrac a4$ if $I=\dfrac{mg}{4aB_0}$}
\end{qoptions}
\qn{3} In a circular region centered at O and having radius $R$, magnetic field increases at a constant rate $a_0$ T/sec. A wire frame OPQ having resistance per unit length $\lambda\ \Omega$/m is placed as shown in the figure. Choose the CORRECT option(s).
\qfig{p2_phy_s1_q03_1}
\begin{qoptions}
\opt{A}{$V_P=V_0$ for any $\theta$ as there is no induced emf across OP.}
\opt{B}{$V_O-V_Q=\dfrac{a_0\pi R^2}{8\left(\sqrt2+1\right)}$ for $\theta=45^\circ$}
\opt{C}{$V_O-V_Q=\dfrac{a_0\pi R^2}{2\left(\sqrt2-1\right)}$ for $\theta=45^\circ$}
\opt{D}{$V_P-V_Q=\dfrac{a_0\pi R^2}{8\sqrt2}$ for $\theta=45^\circ$}
\end{qoptions}
\qn{4} In the figure shown, the inductors do not carry any current at the initial moment. The switch is closed at $t=0$. There is no resistance in the circuit. Choose the correct statement (s) :
\qfig{p2_phy_s1_q04_1}
\begin{qoptions}
\opt{A}{The maximum current through the inductors is at $t=\dfrac\pi2\sqrt{\dfrac{2LC}{3}}$.}
\opt{B}{The charge on left plate becomes $-Q_0/2$ at $t=\dfrac{5\pi}{6}\sqrt{\dfrac{3LC}{2}}$.}
\opt{C}{The maximum current through the inductor of inductance $L$ is $Q_0\sqrt{\dfrac{2}{3LC}}$.}
\opt{D}{The maximum rate of change of current through the inductor of inductance $2L$ is $\dfrac{Q_0}{2LC}$.}
\end{qoptions}
\qn{5} A rigid uniform stick of length $L$ is restricted to move on the horizontal X-Y plane. A point charge $q$ is fixed on one end of the stick and a point charge $-q$ is fixed on the other end. At a particular moment the velocity of the stick center is $\vec v=v_x\hat i+v_y\hat j$, the stick is at an angle $\theta$ ($<90^\circ$) to the X-axis, and is spinning around its center at angular velocity $\omega$ in anticlockwise sense. A uniform magnetic field $B$ is applied along the Z-axis.
\begin{qoptions}
\opt{A}{The net force acting on the stick is zero.}
\opt{B}{The net force acting on the stick is $+qBL\omega\left(\cos\theta\,\hat i+\sin\theta\,\hat j\right)$.}
\opt{C}{The net torque acting on the stick is zero.}
\opt{D}{The net torque acting on the stick is non-zero.}
\end{qoptions}
\qn{6} A cylindrical thin shell has length $\ell$ and radius $a$, where $\ell\gg a$. The surface charge density on the shell is $\sigma$. The shell rotates about its axis with an angular velocity $\omega$ which increases slowly with time as $\omega=kt$, where $k$ is a constant.
\qfig{p2_phy_s1_q06_1}
\begin{qoptions}
\opt{A}{The magnetic field inside the cylinder is $\dfrac{\mu_0\sigma\omega a}{2}$.}
\opt{B}{The electric field inside the cylinder is $\dfrac{\mu_0\sigma akr}{2}$.}
\opt{C}{The magnetic energy inside the cylinder is $\dfrac{\pi\mu_0\sigma^2a^4k^2\ell t^2}{2}$.}
\opt{D}{The electric energy inside the cylinder is $\dfrac{\pi\epsilon_0\mu_0^2\sigma^2k^2a^6\ell}{16}$.}
\end{qoptions}
\sect{SECTION-II}
\qn{1} Consider a smooth conducting ring in horizontal plane in region of vertical uniform magnetic field, $B=1$T. External circuit is attached at center and periphery of conducting ring which contain battery of EMF 4 volt and inductance $L=1$H. Circuit is closed at $t=0$, smooth conducting movable rod (mass 1 kg, length 1m) attaches center and periphery of ring. Neglect magnetic field produced by circuit, find the maximum angular velocity of rod (in rad/sec) in subsequent motion.
\qfig{p2_phy_s2_q01_1}
\qn{2} In an equipment, two 0.5 m long conducting rods are rotating at the same speed in opposite directions, and both are perpendicular to a 4 T magnetic field. As the drawing shows, the ends of these rods come to within 1.0 mm of each other as they rotate. Moreover, the fixed ends about which the rods are rotating are connected by a wire, so these ends are at the same electric potential. It is known that the breakdown strength of air is $3\times10^6$ V/m. The right rod is rotating in the clockwise direction and the left rod in the anticlockwise direction. What is the angular speed (in rad/s) of the rods when a spark jumps across the gap?
\qfig{p2_phy_s2_q02_1}
\qn{3} The following diagram depicts a single wire that is bent into the shape below. The circuit is placed in a magnetic field pointing out of the page, uniformly increasing at the rate $\dfrac{dB}{dt}=2.34$ T/s. Calculate the magnitude of induced electromotive force (in Tm$^2$s$^{-1}$) in the wire, in terms of the following labelled areas (m$^2$). Note that B is non-inclusive of C and that A $=4.23$, B $=2.74$, and C $=0.34$.
\qfig{p2_phy_s2_q03_1}
\qn{4} A coil with 5000 turns, $r_1=1.25$ cm and $\ell_1=1.0$ m has a core with $\mu_r=50$. A second coil of 500 turns, $r_2=2.0$ cm and $\ell_2=10.0$ cm is concentric with the first coil. The space between the coils is vaccum. The mutual inductance (in mH) is :
\qn{5} The uniform rod of mass $m$ is freely pivoted about a horizontal axis through point O. A spring having spring constant $k$ and a damper having damping constant $c$ is attached to the rod as shown in the diagram. The force due to damper is equal to $-c\vec v$. Assuming small oscillations, the value of $c$ (in N-s/m) for which the oscillations are critically damped is (Take : $a=1$m, $b=2$m, $k=100$ N/m, $m=3$ kg) (Neglect gravity)
\qfig{p2_phy_s2_q05_1}
\qn{6} A point-like negatively charged particle is projected from a point P of a uniform electric field having magnitude 100 N/C in the direction perpendicular to the electric field with a velocity $v_0$. Uniform magnetic field having magnitude 10 T is also present, which is perpendicular to both the electric field vector $E$ and to the velocity $v_0$. The two types of fields are separated by a plane, which is perpendicular to the electric field vector, as shown in figure. (The whole arrangement is in vacuum, and the effect of the gravitational force on the particle is negligible). The value of $v_0$ for which the charge particle returns to its initial point of projection is
\qfig{p2_phy_s2_q06_1}
\qfullbegin
\sect{SECTION-III}
\qn{1} As seen from figure, currents $I_1$ and $I_2$ can be independently increased, decreased, or held constant by varying resistances $R_1$ and $R_2$. Let $L_1=50$ mH, $L_2=40$ mH, $M=15$ mH. $I_1$ is increased at the rate of 120 A/s and $I_2$ is decreased at the rate of 200 A/s. Calculate the magnitude of induced emf in coil-1 (in V).
\qfig{p2_phy_s3_q01_1}
\qfullend
\qn{2} The coil consists of a middle cylinder with a radius $r$ and two extreme cylinders with radii $R>r$. A long thin wire is tightly wound onto a reel as follows: first one of the outer cylinders is wound, and then the same wire is continued to be wound onto the middle cylinder in the same direction in which the winding began. After winding is completed, the coil is placed on a horizontal table placed in a uniform constant magnetic field $B$, whose lines of induction are parallel to the axis of the coil. An ideal voltmeter is connected to the first end of the wire lying on the table, and the other end of the wire, touching the fixed sliding contact connected to the voltmeter, is pulled along the surface of the table at a constant speed $v$ in a direction perpendicular to the coil axis (see figure). Assuming that the coil rolls on the table without slipping, find the voltmeter reading (in V). Take : $R=2$m, $r=1$m, $v=2$m/s, $B=9$T.
\qfig{p2_phy_s3_q02_1}
\qn{3} A thin non conducting disc of mass $M=2$kg, charge $Q=2\times10^{-2}$C and radius $R=\dfrac16$m is placed on a frictionless horizontal plane with its centre at the origin of the coordinate system. A non uniform, radial magnetic field $\vec B=B_0\hat r$ exists in space, where $B_0=10$T and $\hat r$ is a unit vector in the radially outward direction. The disc is set in motion with an angular velocity $\omega=x\times10^2$ rad/sec, about an axis passing through its centre and perpendicular to its plane, as shown in figure. At what value of $x$, the disc will lift off from the surface.
\qfig{p2_phy_s3_q03_1}
\qn{4} The equivalent inductance (in H) between points P and Q in the figure is :
\qfig{p2_phy_s3_q04_1}
\qn{5} A rectangular loop (10 cm $\times$ 15 cm) lies in X-Y plane with the longer side on $x$-axis and shorter side on $y$-axis and one end of the rectangle at the origin. A time dependent magnetic field $\vec B=B_0(x\sin\omega t\,\hat j+y\cos\omega t\,\hat k)$ with $B_0=120$ gauss/m and $\omega=1000$ rad/s exists within the loop. Amplitude of induced emf (in mV) is
\qn{6} A linear harmonic oscillator having a mass of 1.10 kg is set into motion with viscous damping. If the frequency is 10 Hz and if two successive amplitudes a full cycle apart are measured to be 4.65 mm and 4.30 mm, the viscous damping coefficient (in N-s/m) is (in nearest integer). For $x\ll1$, $e^x=(1+x)$.
\partsec{2}{CHEMISTRY}
\sect{SECTION-I}
\qn{1} Which of the following combination of aqueous solutions of equal concentrations can result in the formation of a buffer solution ?
\begin{qoptions}
\opt{A}{HCl(aq) + NH$_4$OH(aq)}
\opt{B}{NaOH(aq) + C$_6$H$_5$COOH(aq)}
\opt{C}{NH$_4$OH(aq) + HCN(aq)}
\opt{D}{HClO$_4$(aq) + KOH(aq)}
\end{qoptions}
\qn{2} Which of the following properties of a system is/are extensive ?
\begin{qoptions}
\opt{A}{Standard emf of the cell}
\opt{B}{pH of an acidic solution}
\opt{C}{Internal energy}
\opt{D}{Gibb's free energy.}
\end{qoptions}
\qn{3} Following data is given at 300 K :
\begin{center}
\setlength{\tabcolsep}{3pt}\renewcommand{\arraystretch}{1.3}\footnotesize
\begin{tabular}{|l|c|c|}\hline
 & $\Delta H^0_f$ (kJ mol$^{-1}$) & $S^0_m$ (JK$^{-1}$mol$^{-1}$)\\\hline
CO$_2$(g) & $-393$ & 214\\\hline
CO(g) & $-110$ & 200\\\hline
C(graphite) & 0 & 6\\\hline
\end{tabular}
\end{center}
For the reaction\par
CO$_2$(g) + C(s) $\longrightarrow$ 2CO(g) at 300 K \& 1 bar
\begin{qoptions}
\opt{A}{$\Delta_rH^0=173$ kJ/ mol}
\opt{B}{$\Delta_rS^0=180$ Jmol$^{-1}$K$^{-1}$}
\opt{C}{$\Delta_rG^0=53.8$ kJ/ mol}
\opt{D}{$K_p^0<1$}
\end{qoptions}
\qn{4} The salt Zn(OH)$_2$ is involved in the following two equilibria,\par
Zn(OH)$_2$(s) $\rightleftharpoons$ Zn$^{2+}$(aq) + 2OH$^-$(aq) \quad ; $K_{sp}=10^{-17}$\par
Zn(OH)$_2$(s) + 2OH$^-$(aq) $\rightleftharpoons$ [Zn(OH)$_4$]$^{2-}$(aq.) \ ; $K_c=0.1$\par
For minimum solubility of Zn(OH)$_2$(s) :
\begin{qoptions}
\opt{A}{[H$^+$] $=10^{-10}$ M}
\opt{B}{[H$^+$] $=10^{-4}$ M}
\opt{C}{[OH$^-$] $=10^{-4}$ M}
\opt{D}{pH $=7$}
\end{qoptions}
\qn{5} For the hypothetical reaction\par
A$_2$(g) + B$_2$(g) $\rightleftharpoons$ 2AB(g)\par
$\Delta_rG^\circ$ and $\Delta_rS^\circ$ are 20 kJ/mol and $-20$ JK$^{-1}$ mol$^{-1}$ respectively at 200 K.\par
Select the correct option(s). (Use ln 2 $=0.7$)\par
[Given : $C_P$(AB, g) $=40$ JK$^{-1}$ mol$^{-1}$, $C_P$(A$_2$, g) $=25$ JK$^{-1}$ mol$^{-1}$, $C_P$(B$_2$, g) $=35$ JK$^{-1}$ mol$^{-1}$]
\begin{qoptions}
\opt{A}{$\Delta G^\circ$ at 400 K $=22.4$ kJ mol$^{-1}$.}
\opt{B}{$\Delta H^\circ$ at 400 K $=20$ KJ mol$^{-1}$.}
\opt{C}{$\Delta S^\circ$ at 400 K $=-6$ JK$^{-1}$ mol$^{-1}$.}
\opt{D}{$\Delta H^\circ$ at 200 K $=-16$ KJ mol$^{-1}$.}
\end{qoptions}
\qn{6} pH of 0.1 M solution of NaA is 8.92. If a drop of HPh ($pK_a=9.52$) is added to the above solution, then :\par
(It is the medical fact that our eyes can see pink colour if the molar \% of ionized form of indicator is 25\% or more.)\par
(log 2 $=0.30$)
\begin{qoptions}
\opt{A}{$pK_a$ of HA $=4.74$}
\opt{B}{Pink colour will be visible}
\opt{C}{Pink colour will not be visible}
\opt{D}{$pK_a$ of HA $=4.84$}
\end{qoptions}
\sect{SECTION-II}
\qn{1} Potassium hydrogen tartrate, KHC$_4$H$_4$O$_6$, is sometimes used in the detection of potassium. For tartaric acid, the first ionisation constant is $1.0\times10^{-3}$. What is the least molarity of K$^+$ ion which must be present in 1.0 M $-$ H$_2$C$_4$H$_4$O$_6$ (tartaric acid) solution for a precipitate to form ? The solubility product of potassium hydrogen tartrate is $4.74\times10^{-4}$. (Ignore the effect of second dissociation). If the answer is $x\times10^{-2}$ M, the value of `x' is :\par
[Use $\sqrt{10}=3.16$]
\qn{2} How many milligram of sodium acetate should be added in 500 ml of 0.2 M acetic acid solution in order to make the [H$^+$] in the solution, $4\times10^{-4}$ M. $K_a$ of CH$_3$COOH $=1.8\times10^{-5}$. [Atomic mass : Na $=23$]
\qn{3} Calculate magnitude of work involved (in J) when 100 gm of Calcite form of CaCO$_3$ is converted into its Aragonite form of CaCO$_3$ at 100 bar pressure.\par
Given density of Calcite $=2$gm/cc and density of Aragonite $=2.5$ gm/cc.\par
[Atomic mass : Ca $=40$]
\qn{4} An amount of 0.01 moles of solid AgCN is rendered soluble in one litre of aqueous solution by adding just sufficient excess cyanide ion to form complex ion Ag(CN)$_2^-$ and the concentration of free cyanide ion at equilibrium is $2.5\times10^{-7}$ M. Determine [Ag$^+$] $\times10^{10}$ in the solution neglecting hydrolysis of cyanide ion. $K_{diss}$ for Ag(CN)$_2^-$ $=1.0\times10^{-20}$.
\qn{5} One mole of a monoatomic ideal gas initially at a pressure of 2.00 bar and a temperature of 273 K is taken to a final pressure of 4.00 bar by a reversible path defined by equation : $P(V)^{-3/2}=$ constant , the value of $\left|\dfrac{\Delta U}{w}\times100\right|$ for this process is?
\qn{6} Consider the following reversible reaction,\par
A(g) + B(g) $\rightleftharpoons$ C(g)\par
The activation energy of the backward reaction exceeds that of the forward reaction by $RT$ (in J mol$^{-1}$). If the pre-exponential factor of the forward reaction is 2 times that of the reverse reaction the absolute value of $\Delta G^\theta$ (in J mol$^{-1}$) for the reaction at 300 K is.\par
(Given : ln (2) $=0.7$, $RT=2500$ J mol$^{-1}$ at 300 K and $G$ is the Gibbs energy)
\sect{SECTION-III}
\qn{1} 1500 ml 0.1 M H$_3$PO$_4$ is added to 500 mL 0.1 M Na$_3$PO$_4$. Find pH of the final solution.\par
Given, for H$_3$PO$_4$ : $pK_{a_1}=2.5$, $pK_{a_2}=7.5$ and $pK_{a_3}=12.5$ at the given temperature.\par
Use : log2 $=0.3$, log3 $=0.5$, log5 $=0.7$
\qn{2} A reversible cyclic process involves 6 steps. In steps 1 and 3 system absorbs 500 J and 800 J of heat respectively from a heat reservoirs at constant temperatures of 250K and 200K respectively. Steps 2, 4 and 6, are adiabatic such that the temperature of one reservoir changes to that of next. Total work done by the system in a complete cycle is 700 J. Find the value of $\dfrac{T_5}{100}$ if during step 5 system exchanges heat with a reservoir at constant temperature $T_5$(K).
\qn{3} An ideal gas is taken in a thermally insulated container and expanded again vacuum. Calculate the number of parameters/quantities out of the following which will be zero for this process :\par
$\Delta U$, $\Delta H$, q, W, $\Delta S_{system}$, $\Delta S_{surrounding}$, $\Delta V$, $\Delta P$, $\Delta G$, $\Delta T$\par
(Assuming that gas is not showing association or dissociation)
\qn{4} $K_a$ of HCN is $5\times10^{-10}$ at 25$^\circ$C. For maintaining a constant pH of 9, the volume (in mL) of 5M KCN solution required to be added to 10mL of 2M HCN solution is :
\qn{5} A sample of an ideal monoatomic gas undergoes a reversible process following equation :\par
$T^3\cdot V^x=$ constant.
\begin{center}
$\left[\begin{array}{c}\text{State}-\text{I}\\ T_1,V_1,P_1\end{array}\right]\longrightarrow\left[\begin{array}{c}\text{State}-\text{II}\\ T_2,V_2,P_2\end{array}\right]$
\end{center}
If $T_2=4\times T_1$ and entropy of gas at both states I and II are equal, then the value of `x' is :
\qn{6} Silver ion complexes with cyanide ion to form Ag(CN)$_2^-$ with formation constant, \ $K_f=10^{22}$. On adding $5\times10^{-3}$ moles of silver nitrate to 1 litre of 0.01 M NaCN solution if final concentration of Ag$^+$ ions in solution is [P $\times10^{-Q}$] then find value of [Q $-$ P] ?
\partsec{3}{MATHEMATICS}
\sect{SECTION-I}
\qn{1} Let A$(4,-2)$ and B$(0,1)$ are two points on a parabola. Normals at A \& B to the parabola intersect at P$(2,2)$. Tangent parallel to AB to the parabola touches the parabola at Q and tangents drawn at A \& B to the parabola meet at R.
\begin{qoptions}
\opt{A}{Equation of directrix of the parabola is $2x+y+1=0$}
\opt{B}{Equation of directrix of the parabola is $y=-3$}
\opt{C}{Length of latus rectum of parabola is $\dfrac85$}
\opt{D}{Area of triangle AQB is equal to $\dfrac52$}
\end{qoptions}
\qn{2} Let tangent at point P on ellipse cuts directrix corresponding to focus $S_1$ and $S_2$ of ellipse at $D_1$ and $D_2$ respectively. This tangent cuts auxiliary circle at $F_1$ and $F_2$ (where $F_1D_1<F_2D_1$). Lines $D_1S_1$ and $D_2S_2$ meet at Q. If $F_1F_2=$ length of semi major axis then
\begin{qoptions}
\opt{A}{$D_1Q=D_2Q$}
\opt{B}{$D_1Q\neq D_2Q$}
\opt{C}{$\angle PD_1S_1=30^\circ$}
\opt{D}{$\angle PD_2S_2=60^\circ$}
\end{qoptions}
\qn{3} Let $e$ be the eccentricity of a hyperbola $\dfrac{x^2}{a^2}-\dfrac{y^2}{b^2}=1$. If $f(e)$ is the eccentricity of the conjugate hyperbola of hyperbola whose eccentricity is $e$, then $\displaystyle\int_1^3\underbrace{f(f(f(\ldots f(e)\ldots)))}_{n\text{ times}}\,de$ is equal to
\begin{qoptions}
\opt{A}{2, if $n$ is even}
\opt{B}{4, if $n$ is even}
\opt{C}{$2\sqrt2$, if $n$ is odd}
\opt{D}{$4\sqrt2$, if $n$ is odd}
\end{qoptions}
\qn{4} Let A$(at_1^2,2at_1)$, B$(at_2^2,2at_2)$, C$(at_3^2,2at_3)$, be three points on the parabola $y^2=4ax$. If the orthocentre of triangle ABC is focus S of the parabola, then
\begin{qoptions}
\opt{A}{$t_1t_2+t_3t_2+t_3t_1=-5$}
\opt{B}{$(1+t_1)(1+t_2)(1+t_3)=-4$}
\opt{C}{If $t_1=0$ then $t_2+t_3=0$}
\opt{D}{$\dfrac1{t_1t_2}+\dfrac1{t_2t_3}+\dfrac1{t_3t_1}=1$}
\end{qoptions}
\qn{5} Choose the correct statement(s) :
\begin{qoptions}
\opt{A}{A hyperbola whose asymptotes are $x+2y+3=0$ and $3x+4y+5=0$, is passing through the point $(1,-1)$. Then the equation of its conjugate hyperbola is $3x^2+10xy+8y^2+\lambda x+22y+\mu=0$ then $\mu-\lambda=6$}
\opt{B}{The minimum and maximum distance of a point $\left(\dfrac95,\dfrac{12}5\right)$ from the ellipse $4(3x+4y)^2+9(4x-3y)^2=900$ are $\lambda$ and $\mu$ respectively then $24\mu-7\lambda=144$}
\opt{C}{Line $x+y=7$ touches a parabola at P(3, 4). If the focus of the parabola is (1, 1) and its vertex is $(\lambda,\mu)$ then $\lambda\mu=12$}
\opt{D}{If the length of the latus rectum of the parabola $169\{(x-1)^2+(y-3)^2\}=(5x-12y+19)^2$ is $\dfrac\lambda\mu$ (where $\lambda$ and $\mu$ are relatively primes) then $\mu+\lambda=37$}
\end{qoptions}
\qn{6} Consider a parabola $y^2=\ell x$ ($\ell>0$) with its vertex at the origin `O' then which of the following statement(s) is(are) true ?
\begin{qoptions}
\opt{A}{If origin is joined to any point P on the given parabola and a line perpendicular to OP through P meet the axis at Q, then projection of PQ on the axis is equal to `$\ell$'}
\opt{B}{A normal is drawn at any arbitrary point P on the parabola which meets the axis at G. Locus of the middle point of PG is a parabola whose latus rectum is $\dfrac\ell4$}
\opt{C}{A variable parabola of latus ractum $\ell$, touches the given parabola, the axes of the two parabolas being parallel. The locus of the vertex of the variable parabola is another parabola whose latus rectum is `$2\ell$'}
\opt{D}{The locus of the mid point of the focal radii of a variable point P moving on the given parabola, is a parabola whose latus rectum is `$4\ell$'}
\end{qoptions}
\sect{SECTION-II}
\qn{1} Let an ellipse $E_1:\dfrac{x^2}{64}+\dfrac{y^2}{48}=1$ and two parabolas $P_1$ and $P_2$ both passing through ends of minor axis of $E_1$. Vertex of each of two parabolas lie on major axis of $E_1$. If $x-2y+16=0$ is common tangent of all the three curves and A be the area bounded by $P_1$ and $P_2$ then $\dfrac{A}{\sqrt3}$ is
\qn{2} An ellipse is inscribed in the circle $x^2+y^2=9$ whose major axis is $3x-4y=0$. From Point $(0,3)$ a perpendicular is drawn to the line $3x-4y=0$ intersecting the ellipse at P and line at Q such that $PQ=1$. If $d$ is the distance between both the foci of ellipse then $\dfrac{288\,d^2}{119}$ is
\qn{3} If $(\alpha,\beta)$ is the centroid of a triangle PQR where P, Q and R be three points lying on the parabola $(y-4)^2=15(x+3)$ and normals at P and Q meet at R. Then the largest integral value of $m$ satisfying $\alpha+\beta>m$ for all possible values of $\alpha$ and $\beta$, is
\qn{4} Let $\vec a=(x^2+3)\hat i-3y^2\hat j+\hat k$ and $\vec b=\hat i-3\hat j-4x\hat k$ be perpendicular vectors for $x,y\in R$ and the relation obtained is $x$ and $y$ represent a conic C. Let $m$ and $n$ are number of integers in the range of values of $x$ and $y$ respectively and if P$(x,y)$ lies on conic C then the equation of director circle of the circle which is tangent to conic C at M(3, 0) and passing through the centre of conic is $ax^2+by^2-cx+23=0$ then the value of $(m^2+n^2+a+b+c)$ is
\qn{5} Consider the following statements :\par
\textbf{Statement-I :} If length of the common tangent of circle $x^2+y^2=4$ and conic $3x^2+5y^2=15$ is $\ell$ and $A=\dfrac1\ell$.\par
\textbf{Statement-II :} The number of possible tangents which can be drawn to the curve $4x^2-9y^2=36$, which are perpendicular to the straight line $5x+2y-10=0$ is B.\par
\textbf{Statement-III :} With one focus of the hyperbola $\dfrac{x^2}{9}-\dfrac{y^2}{16}=1$ as the centre, a circle is drawn which is tangent to the hyperbola with no part of the circle being outside the hyperbola. The radius of the circle is C.\par
\textbf{Statement-IV :} For each positive integer $n$, consider the point P with abscissa $n$ on the curve $y^2-x^2=1$. If $d_n$ represents the shortest distance from the point P to the line $y=x$ then $\displaystyle\lim_{n\to\infty}(n\cdot d_n)$ has the value equal to D.\par
Then the value of $(A+B+C)\div D^2$ is
\qn{6} Let A and B be two points on the major axis of the ellipse $\dfrac{x^2}{25}+\dfrac{y^2}{16}=1$, which are equidistant from the centre. If C and D are the images of these points in the line mirror $y=mx$, $m\neq0$ then the maximum area of quadrilateral ACBD is
\sect{SECTION-III}
\qn{1} An isosceles $\Delta$ABC is inscribed in the circle $x^2+y^2=16$. From points A, B and C three ordinates are drawn which cut the ellipse $\dfrac{x^2}{16}+\dfrac{y^2}{9}=1$ at the points P, Q and R respectively such that A and P on the opposite sides, Q and B are on the same side and C and R are on the same side of the major axis. If $\Delta$PQR is right angled at P and $\theta$ is the smallest angle of the $\Delta$ABC, then the value of $(4\tan\theta+3)$ is
\qn{2} If tangent drawn to the parabola $y^2=-ax$ ($a>0$) from a point A(1, 0) which also touches the hyperbola $\dfrac{x^2}{4}-\dfrac{y^2}{b^2}=1$ at B such that $\angle ASB=90^\circ$, where S is the focus of the hyperbola then the value of $(a-b^2)$ is
\qn{3} Tangents are drawn from any point on straight line $y=x+8$ to the auxiliary circle of $x^2+16y^2=16$. If A and B are the points of contact of these tangent and P, Q are corresponding points of A and B on the ellipse respectively. If locus of mid point of PQ is $x^2+\alpha y^2+\beta x+\gamma y=0$ then the value of $(\alpha-\beta+\gamma)$ is :-
\qn{4} A ray of light emerging from the point P(0, 4) along the tangent of the ellipse $\dfrac{x^2}{25}+\dfrac{y^2}{16}=1$ and strike a circle of radius $\sqrt{65}$ units and reflected from it. The reflected ray entering inside the ellipse through the focus $(-3,0)$ and strike lower surface of ellipse and gets reflected from it. If the reflected ray returns to the point P(0, 4), then ordinate of the centre of this circle is
\qn{5} From a point common tangents are drawn to the circle $x^2+y^2=\dfrac{a^2}{2}$ and parabola $y^2=4ax$. If the area of the quadrilateral formed by the common tangents, the chord of contact of the circle and the chord of contact of the parabola is $\dfrac{Aa^2}{B}$ (where A and B are coprime) then the value of $(A-2B)$ is
\qn{6} A triangle is circumscribed to the ellipse $\dfrac{x^2}{a^2}+\dfrac{y^2}{b^2}=1$ and two of the vertices lie on the directrices and if the locus of the third vertex is the ellipse $b^2y^2=a^2(e^2+\lambda)(a^2-\mu x^2)$ then the value of $(\lambda+\mu)$ is

\end{multicols}
\end{document}
